\section{Method}

In this section, we introduce our method for high-resolution text-to-long video generation.
We first generate $256 \times 256$ resolution long videos (240 frames, or 1200 frames), then enhance them to higher resolution ($720\times 720$).
The overview of the whole pipeline is provided in \figrefcustom{fig:high_level_pipeline}.
The long video generation process comprises three stages: the \textbf{\textit{Initialization Stage}}, where the first 16-frame chunk is synthesized by a pre-trained text-to-video model (\eg Modelscope \cite{Modelscope}), the \textbf{\textit{Streaming T2V Stage}} where new content for subsequent frames is generated autoregressively. To ensure seamless transitions between chunks, we introduce (see \figrefcustom{fig:pipeline}) our conditional attention module (CAM), which utilizes short-term information from the last $F_{cond}=8$ frames and our appearance preservation module (APM), which extracts long-term information from an  anchor frame to maintain object appearance and scene details during the autoregressive process. 
%
After generating a long video (\eg 240, 1200 frames or more), the \textbf{\textit{Streaming Refinement Stage}} enhances the video using a high-resolution text-to-short-video model (\eg MS-Vid2Vid-XL \cite{2023i2vgenxl}) autoregressively with our randomized blending approach for seamless chunk processing. This step does not require additional training, making our approach cost-effective. 


\subsection{Conditional Attention Module}
\label{ssec:cam}
For training a conditional network in our Streaming T2V stage, we leverage the capabilities of a pre-trained text-to-video model (\eg Modelscope \cite{Modelscope}) as a prior for autoregressive long video generation. Subsequently, we will denote this pre-trained text-to-(short)video model as \textbf{\textit{Video-LDM}}.
%
To condition Video-LDM autoregressively with short-term information from the preceding chunk (see \figrefcustom{fig:high_level_pipeline}, mid), we introduce the \textit{\textbf{Conditional Attention Module (CAM)}}. CAM consists of a feature extractor and a feature injector into the Video-LDM UNet, inspired by ControlNet \cite{controlnet}.
%
The feature extractor utilizes a frame-wise image encoder  $\mathcal{E}_{\text{cond}}$, followed by the same encoder layers that the Video-LDM UNet uses up to its middle layer (initialized with the UNet's weights).
For the feature injection, we let each long-range skip connection in the UNet \textit{attend} to corresponding features generated by CAM via cross-attention. 

Let $x$ denote the output of $\mathcal{E}_{\text{cond}}$ after zero-convolution. 
We use addition to fuse $x$ with the output of the first temporal transformer block of CAM. 
For the injection of CAM's features into the Video-LDM Unet, we consider the UNet's skip-connection features $x_{\mathrm{SC}} \in \mathbb{R}^{b \times F \times h \times w \times c}$ (see \figrefcustom{fig:pipeline}) with batch size $b$.  We apply spatio-temporal group norm, and a linear map $P_{\mathrm{in}}$ on $x_{\mathrm{SC}}$. 
Let $x'_{\mathrm{SC}} \in \mathbb{R}^{(b\cdot w \cdot h) \times F \times c}$ be the resulting tensor after reshaping. 
We condition $x'_{\mathrm{SC}}$ on the corresponding CAM feature $x_{\mathrm{CAM}} \in \mathbb{R}^{(b\cdot w \cdot h) \times F_{\mathrm{cond}} \times c}$ (see \figrefcustom{fig:pipeline}), where $F_{\mathrm{cond}}$ is the number of conditioning frames, via temporal multi-head attention (T-MHA) \cite{vaswani2017attention}, \ie independently for each spatial position (and batch).
Using learnable linear maps $P_Q,P_K,P_V$, for queries, keys, and values, we apply T-MHA using keys and values from $x_{\mathrm{CAM}}$ and queries from $x'_{\mathrm{SC}}$, \ie with $Q = P_Q(x'_{\text{SC}}), K = P_K(x_{\mathrm{CAM}}), V =P_V(x_{\mathrm{CAM}})$,
\begin{equation}
    %\begin{split}
       \displaystyle x''_{\mathrm{SC}} = \text{T-MHA}\bigl(Q, K, V )\bigr.
    %\end{split}
\end{equation}

Finally, we use a linear map $P_{out}$ and a reshaping operation $R$, the  output of CAM is added to the skip connection (as in ControlNet \cite{controlnet}): 
\begin{equation}
    x'''_{\mathrm{SC}} = x_{\mathrm{SC}} + R(P_{\text{out}}(x''_{\text{SC}})),
\end{equation}
and $x_{\mathrm{SC}}'''$ is used in the decoder layers of the UNet. 
We zero-initialize $P_{\text{out}}$, so that CAM initially does not affect the base model's output, which improves convergence. 

%CAM utilizes the last $F_{\mathrm{cond}}$ conditional frames of the preceding chunk as input. 
The design of CAM enables conditioning the $F$ frames of the base model on the $F_{\mathrm{cond}}$ frames of the preceding chunk.
In contrast, sparse encoder \cite{guo2023sparsectrl} employs convolution for feature injection, thus needs additional $F-F_{\mathrm{cond}}$ zero-valued frames (and a mask) as input, in order to add the output to the $F$ frames of the base model. These  inconsistencies in the input lead to severe inconsistencies in the output (see 
%% EXTERNAL REFERENCE ATTENTION
% \secrefcustom{sec:appendix.ablation_studies.cam} 
Sec. \textcolor{linkblue}{A.3.1}
and \secrefcustom{sec:exp_comparison_with_baselines}).
\subsection{Appearance Preservation Module}
\label{ssec:apm}
Autoregressive video generators typically suffer from forgetting initial object and  scene features, leading to severe appearance changes. 
To tackle this issue, we incorporate long-term memory by leveraging the information contained in a fixed anchor frame of the very first chunk using our proposed \textbf{Appearance Preservation Module (APM)}.  This helps to maintain scene and object features across video chunk generations
%% EXTERNAL REFERENCE ATTENTION
% (see \figrefcustom{fig:ablation}).
(see Fig. \textcolor{linkblue}{A.8}).

To enable APM to balance guidance from the anchor frame and the text instructions, we propose (see \figrefcustom{fig:pipeline}): (i) We combine the CLIP \cite{CLIP_paper} image token of the anchor frame with the CLIP text tokens from the textual instruction by expanding the clip image token to $k=16$ tokens using an MLP layer, concatenating the text and image encodings at the token dimension, and utilizing a projection block, leading to $x_{\mathrm{mixed}} \in \mathbb{R}^{b \times 77 \times 1024}$; (ii) For each cross-attention layer $l$, we introduce a weight $\alpha_{l} \in \mathbb{R}$ (initialized as $0$) to perform cross-attention using keys and values derived from a weighted sum $x_{\mathrm{mixed}}$, and the usual CLIP text encoding of the text instructions $x_{\mathrm{text}}$: 
\begin{equation}
    x_{\mathrm{cross}} = \mathrm{SiLU}(\alpha_{l}) x_{\mathrm{mixed}} + x_{\mathrm{text}}.
\end{equation}


%To enable APM to balance the guidance by the  anchor frame with the guidance by the text instructions, we propose (see \figrefcustom{fig:pipeline}): (i) We mix the CLIP \cite{CLIP_paper} image token of the anchor frame with the CLIP text tokens from the textual instruction by expanding the clip image token to $k=8$ tokens using a linear layer, concatenate the text and image encodings at the token dimension, and use a projection block, leading to $x_{\mathrm{mixed}} \in \mathbb{R}^{b \times 77 \times 1024}$; (ii) We introduce for each cross-attention layer $l$ a weight $\alpha_{l} \in \mathbb{R}$ (initialized as $0$) to perform cross-attention  using keys and values coming from a weighted sum of $x_{\mathrm{mixed}}$ and the usual CLIP text encoding of the text instructions $x_{\mathrm{text}}$: 
%\begin{equation}
%    x_{\mathrm{cross}} = \mathrm{SiLU}(\alpha_{l}) x_{\mathrm{mixed}} + x_{\mathrm{text}}.
%\end{equation}

The experiments in 
%% EXTERNAL REFERENCE ATTENTION
% \secrefcustom{sec:appendix.ablation_studies.apm} 
Sec. \textcolor{linkblue}{A.3.2}
show that the light-weight APM module helps to keep scene and identity features across the autoregressive process 
%% EXTERNAL REFERENCE ATTENTION
% (see \figrefcustom{fig:ablation}).
(see Fig. \textcolor{linkblue}{A.8}).

\subsection{Auto-regressive Video Enhancement}
\label{ssec:autoreg-enhancer}
To further enhance quality and resolution of our text-to-video results, we use a high-resolution ($1280 \times 720$) text-to-(short)video model (Refiner Video-LDM, see \figrefcustom{fig:high_level_pipeline}), \eg MS-Vid2Vid-XL  \cite{wang2024videocomposer,2023i2vgenxl}, to autoregressively improve 24-frame video chunks. To this end, we add noise to each video chunk and denoise it using Refiner Video-LDM (SDEdit approach \cite{meng2021sdedit}). Specifically, we upscale each low-resolution 24-frame video chunk to $720 \times 720$ using bilinear interpolation \cite{amidror2002scattered}, zero-pad to $1280 \times 720$, encode the frames with the image encoder $\mathcal{E}$ to get a latent code $x_{0}$, apply $T' <T$ forward diffusion steps (see \eqrefcustom{eqn:forward-diffusion}) so that $x_{T'}$  still contains signal information, and denoise it with Refiner Video-LDM.

%To further improve quality and resolution of our text-to-video results, we utilize a high-resolution (1280x720) text-to-(short)video model (Refiner Video-LDM, see \figrefcustom{fig:pipeline}) to autoregressively enhance 24-frame chunks of generated videos. To this end, we add a substantial amount of noise to the input video-chunk and denoise it with the text-to-video diffusion model (SDEdit \cite{meng2021sdedit} approach).
%More precisely, we take a high-resolution text-to-video model (\eg  MS-Vid2Vid-XL \cite{wang2024videocomposer,2023i2vgenxl}), and a low resolution video chunk of 24 frames which is first bilinearly upscaled \cite{amidror2002scattered} to the target high resolution. 
%Then we encode the frames using the image encoder $\mathcal{E}$ so that we obtain a latent code $x_{0}$. 
%After that we apply $T' <T$ forward diffusion steps (see Eq.  \eqref{eqn:forward-diffusion}) so that $x_{T'}$ still contains signal information (mostly about the video structure), and denoise it using the high-resolution video diffusion model.

Naively enhancing each chunk independently leads to inconsistent transitions (see \figrefcustom{fig:video_enhancer_main} (a)).
To overcome this shortcoming, we introduce shared noise and a randomized blending technique.  We divide a low-resolution long video into $m$ chunks $\mathcal{V}_{1},\ldots,\mathcal{V}_{m}$ of $F=24$ frames, each with an $\mathcal{O}=8$ frames overlap between consecutive chunks. For each denoising step, we must sample noise (see \eqrefcustom{eqn:diffusion-backward-process}). We combine that noise with the noise already sampled for the overlapping frames of the preceding chunk to form \textbf{shared noise}. Specifically, for chunk $\mathcal{V}_{i}$, $i=1$, we sample noise $\epsilon_{1} \sim \mathcal{N}(0,I)$ with $\epsilon_{1} \in \mathbb{R}^{F \times h \times w \times c}$. For $i>1$,  we sample noise $\hat{\epsilon}_{i}\sim \mathcal{N}(0,I)$ with $\hat{\epsilon}_{i} \in \mathbb{R}^{(F-\mathcal{O}) \times h \times w \times c}$  and concatenate it with  $\epsilon_{i-1}^{(F-\mathcal{O}):F}$ (already sampled for the preceding chunk) along the frame dimension to obtain $\epsilon_{i}$ \ie:
\begin{equation}
    \epsilon_{i} := \mathrm{concat}([\epsilon_{i-1}^{(F-\mathcal{O}):F},\hat{\epsilon}_{i}],\mathrm{dim}=0).
\end{equation} 
 At diffusion step $t$ (starting from $T'$), we perform one denoising step using $\epsilon_{i}$ and obtain for chunk $\mathcal{V}_{i}$  the latent code $x_{t-1}(i)$. 
 Despite these efforts, transition misalignment persists (see \figrefcustom{fig:video_enhancer_main} (b)).
 
%we must sample noise (see Eq. \ref{eqn:diffusion-backward-process}). Starting from the first chunk $\mathcal{V}_{1}$, we sample noise $\epsilon_{1} \sim \mathcal{N}(0,I)$ with $\epsilon_{1} \in \mathbb{R}^{F \times h \times w \times c}$ and concatenate it with the noise from the overlapping frames of the preceding chunk to create the \textit{shared noise}. After each denoising step, we obtain the latent code for the chunk. Despite these efforts, transition misalignment persists (see \figrefcustom{fig:video_enhancer} (b)).

%To address inconsistent transitions in enhancing video chunks, we introduce shared noise and a randomized blending technique. We divide a low-resolution long video into M chunks of 24 frames each with an 8-frame overlap between consecutive chunks. During backward diffusion, we sample noise for denoising. Starting with the first chunk, we sample noise and concatenate it with the noise from overlapping frames of the previous chunk to create shared noise. After one denoising step, we obtain the latent code for the chunk. Despite these efforts, transition misalignment persists.
%
%However,  the na\"ive approach of independently enhancing each chunk leads to inconsistent transitions (see \figrefcustom{fig:video_enhancer} (a)).
%We tackle this issue by using shared noise between consecutive chunks and leveraging our randomized blending approach.
%Given our low-resolution long video, we split it into $m$ chunks $\mathcal{V}_{1},\ldots,\mathcal{V}_{m}$ of $F=24$ frame-length such that each two consecutive chunks have an overlap of $\mathcal{O}=8$ frames.  For the backward diffusion at step $t$, starting from $T'$, we must sample noise to perform one denoising step (see Eq. \ref{eqn:diffusion-backward-process}). We start with the first chunk $\mathcal{V}_{1}$ and sample noise $\epsilon_{1} \sim \mathcal{N}(0,I)$ with $\epsilon_{1} \in \mathbb{R}^{F \times h \times w \times c}$.  For each subsequent  chunk $\mathcal{V}_{i}$, $i>1$,  we sample noise $\hat{\epsilon}_{i}\sim \mathcal{N}(0,I)$ with $\hat{\epsilon}_{i} \in \mathbb{R}^{(F-\mathcal{O}) \times h \times w \times c}$ and concatenate it along the frame dimension with the noise $\epsilon_{i-1}^{(F-\mathcal{O}):F}$ that was sampled for the $\mathcal{O}$ overlapping frames of the previous chunk, \ie 
%\begin{equation}
%    \epsilon_{i} := \mathrm{concat}([\epsilon_{i-1}^{(F-\mathcal{O}):F},\hat{\epsilon}_{i}],\mathrm{dim}=0), \, \text{for all } i =2,\ldots,m,
%\end{equation} 
%so that we obtain shared noise for overlapping frames. We perform one denoising step using $\epsilon_{i}$ and obtain for chunk $\mathcal{V}_{i}$  the latent code $x_{t-1}(i)$. 
%Yet, this approach is not sufficient to remove transition misalignment (see \figrefcustom{fig:video_enhancer} (b)).
%
To significantly improve consistency, we introduce \textbf{\textit{randomized blending}}.
Consider the latent codes $x_{L} := x_{t-1}(i-1)$ and  $x_{R} := x_{t-1}(i)$ of two consecutive chunks $\mathcal{V}_{i-1},\mathcal{V}_{i}$ at denoising step $t-1$. The latent code $x_{L}$ of chunk $\mathcal{V}_{i-1}$ possesses a smooth transition from its first frames to the overlapping frames, while the latent code $x_{R}$  possesses a smooth transition from the overlapping frames to the subsequent frames. 
Thus, we combine the two latent codes via concatenation at a randomly chosen overlap position, by randomly sampling a frame index $f_{\text{thr}}$ from  $\{0,\ldots,\mathcal{O}\}$ according to which we merge the two latents $x_{L}$ and $x_{R}$:
\begin{equation}
x_{LR} := \mathrm{concat}([x_{L}^{1:F-f_{\text{thr}}},x_{R}^{f_{\text{thr}}+1:F}], \mathrm{dim} = 0).
\end{equation}
%then taking from $x_{t-1}^{(F-\mathcal{O}):F}(i-1)$ the latent code of the first $f_{\text{thr}}$ frames and from  $x_{t-1}^{1:\mathcal{O}}(i)$ the latent code of the frames starting from $f_{\text{thr}}+1$. 
Then, we update the latent code of the entire long video $x_{t-1}$ on the overlapping frames and perform the next denoising step.  Accordingly, for a frame $f \in \{1,\ldots,\mathcal{O} \}$ of the overlap, the latent code of  chunk $\mathcal{V}_{i-1}$ is used with probability $1- \frac{f}{\mathcal{O}+1}$.
%
This probabilistic mixture of latents in overlapping regions effectively diminishes inconsistencies between chunks (see \figrefcustom{fig:video_enhancer_main}(c)).
%
The importance of randomized blending is further assessed in an ablation study in the appendix
%% EXTERNAL REFERENCE ATTENTION
% (see \secrefcustom{sec:appendix.ablation_studies}). 
(see Sec. \textcolor{linkblue}{A.3})


%By using a probabilistic mixture of the latents in an overlapping region, we successfully diminish inconsistencies between chunks (see \figrefcustom{fig:video_enhancer}(c)).

